\begin{theorem}[森田纪一]\label{prop:Morita-refined}
	考虑 $\Bbbk$-代数 $A$ 和 $B$. 命 $\mathrm{Equiv}(\cated{Mod}A, \cated{Mod}B)$ 为以所有等价 $F: \cated{Mod}A \to \cated{Mod}B$ 为对象, 以其间的同构为态射的范畴. 另一方面, 定义范畴 $\mathcal{P}(A, B)$ 使得其对象是满足下述条件的 $(A, B)$-双模 $P$:
	\begin{itemize}
		\item $P$ 作为右 $B$-模是 $\cated{Mod}B$ 的有限生成投射生成元,
		\item 左乘诱导同构 $A \rightiso \End_{\cated{Mod}B}(P)$,
	\end{itemize}
	其态射则定为双模的同构. 我们有以下互为拟逆的函子:
	\[\begin{tikzcd}[row sep=tiny]
		\mathrm{Equiv}(\cated{Mod}A, \cated{Mod}B) \arrow[shift left, r] & \mathcal{P}(A, B) \arrow[shift left, l] \\
		F \arrow[mapsto, r] & F(A) \\
		(\cdot) \dotimes{A} P & P \arrow[mapsto, l] .
	\end{tikzcd}\]
\end{theorem}
\begin{proof}
	首先说明 $F \mapsto F(A)$ 是良定义的. 观察到 $A$ 是 $\cated{Mod}A$ 的有限生成投射生成元, 既然投射生成元和有限生成性质皆有范畴论刻画 (引理 \ref{prop:fg-module-cat}), 故 $F(A)$ 亦然; 此外, 根据抽象的理由,
	\[ F: \End_{\cated{Mod}A}(A) \rightiso \End_{\cated{Mod}B}(F(A)), \]
	不难看出这正是 $A$ 对 $F(A)$ 的左乘给出的 $A \to \End_{\cated{Mod}B}(F(A))$.
	
	若能说明从右到左的函子也是良定义的, 则定理 \ref{prop:Morita} 便蕴涵双向的函子互为拟逆. 问题遂归结为证: 若 $P \in \Obj(\mathcal{P}(A, B))$, 则 $F := (\cdot) \dotimes{A} P$ 是范畴的等价.
	
	取 $P$ 如上. 在 \eqref{eqn:AB-adjunction-2} 已经说明 $(\cdot) \dotimes{A} P$ 的右伴随是 $(\cdot) \dotimes{B} P^\vee$, 其中 $P^\vee := \Hom_{\cated{Mod}B}(P, B)$. 目标是说明它们互为拟逆. 在定义--命题 \ref{def:Morita-equiv} (ii) 中取 $Q = P^\vee$, 将问题进一步化约为证明定义 \ref{def:bimodule-ev-coev} 的双模同态
	\[ \mathrm{ev}: P^\vee \dotimes{A} P \to B, \quad \mathrm{coev}: A \to P \dotimes{B} P^\vee \simeq \End_{\cated{Mod}B}(P) \]
	皆为同构.
	
	先看 $\mathrm{ev}$. 由 $P$ 是 $\cated{Mod}B$ 的生成元可知存在满同态 $P^{\oplus I} \twoheadrightarrow B$, 特别地, 存在 $q_1, \ldots, q_n \in P^\vee$ 和 $p_1, \ldots, p_n \in B$ 使得 $\sum_{i=1}^n q_i(p_i) = 1_B$, 亦即 $\mathrm{ev}(\sum_i q_i \otimes p_i) = 1_B$. 满性得证.
	
	接着引进写作乘法的运算 $P \times P^\vee \to A$, 它的刻画是对于所有 $p, p' \in P$ 和 $q \in P^\vee$, 在 $P$ 中有形似结合律的
	\[ \underbracket{(pq)}_{\in A} p' = p \underbracket{(qp')}_{\in B}, \]
	其中 $qp' := q(p')$; 诚然, 当 $p$ 和 $q$ 固定, 右式是 $\End_{\cated{Mod}B}(P)$ 的元素, 而左乘诱导 $A \rightiso \End_{\cated{Mod}B}(P)$, 由此唯一定义了 $pq \in A$. 从上式还可以推导第二种结合律, 这是 $P^\vee$ 中的等式
	\[ \underbracket{(qp)}_{\in B} q' = q \underbracket{(pq')}_{\in A}; \]
	说明两边对所有 $p' \in P$ 取值皆相等即可, 平凡的验证留给读者.

	现在可以说明 $\mathrm{ev}$ 的单性. 选定上一步得到的 $p_i$ 和 $q_i$, 设 $\mathrm{ev}(\sum_{j=1}^m q'_j \otimes p'_j) = 0$, 则上述两种结合律蕴涵
	\begin{multline*}
		\sum_j q'_j \otimes p'_j = \sum_{i, j} (q_i p_i) q'_j \otimes p'_j
		= \sum_{i, j} q_i \underbracket{(p_i q'_j)}_{\in A} \otimes p'_j \\
		= \sum_{i, j} q_i \otimes (p_i q'_j) p'_j
		= \sum_{i, j} q_i \otimes p_i \underbracket{(q'_j p'_j)}_{\in B}
		= \sum_j q_i \otimes p_i \sum_j q'_j p'_j = 0.
	\end{multline*}

	最后, 关于 $\mathrm{coev}$ 为同构的断言无非是 $\mathcal{P}(A, B)$ 定义的第二部分.
\end{proof}
